% Generated draft: requires mathematical and source review.
\documentclass[11pt]{article}
\usepackage[margin=25mm]{geometry}
\usepackage{fontspec}
\setmainfont{Latin Modern Roman}
\usepackage{amsmath,amssymb,mathtools,unicode-math}
\setmathfont{Latin Modern Math}
\usepackage{enumitem,needspace,xurl,hyperref}
\setlength{\parindent}{0pt}
\setlength{\parskip}{5pt}
\setlength{\emergencystretch}{3em}
\newcommand{\N}{\mathbb{N}}
\newcommand{\Z}{\mathbb{Z}}
\newcommand{\Q}{\mathbb{Q}}
\newcommand{\R}{\mathbb{R}}
\newcommand{\C}{\mathbb{C}}
\newcommand{\F}{\mathbb{F}}
\newcommand{\A}{\mathbb{A}}
\newcommand{\PP}{\mathbb{P}}
\newcommand{\E}{\mathbb{E}}
\newcommand{\eps}{\varepsilon}
\newcommand{\dd}{\,\mathrm{d}}
\begin{document}
\section*{UnsolvedMath ID: 30005627}
\textbf{Source problem code:} catalogue entry
\subsection{Definitions and mathematical statement}
The supplied record concerns an unstable kinetic equilibrium \(f^{\mathrm{eq}}\), perturbed initially as \(f_0=f^{\mathrm{eq}}+\widetilde f\), where \(\widetilde f\) is described as small. It displays the controlled kinetic equation
\[
\begin{cases}
\partial_t f+v\partial_x f-\bigl(H+E[f]\bigr)\partial_v f=0,\\
E[f]=\partial_xG*(1-\rho_f),\\
f(0,x,v)=f^{\mathrm{eq}}(x,v)+\widetilde f(x,v).
\end{cases}
\]
Here \(f=f(t,x,v)\) is the kinetic distribution, \(H\) is an imposed external electric field, and \(E[f]\) is the self-generated electric field. The notation \(*\) denotes spatial convolution. Interpreting \(\rho_f\) as the velocity density gives formally
\[
\rho_f(t,x)=\int_V f(t,x,v)\,\mathrm{d}v,
\]
where the velocity domain \(V\), the spatial domain, the kernel \(G\), and all boundary conditions remain unspecified in the supplied excerpt. The displayed constant \(1\) in the field equation is retained without imposing an unstated normalization or neutrality condition.

\textbf{Problem, subject to unresolved specifications.} For the unstable equilibrium and a small initial perturbation, determine whether an admissible external electric field \(H\) can stabilize the resulting dynamics, and characterize an optimal stabilizing field. The source proposes the formal PDE-constrained optimization problem
\[
\min_H J\bigl(f[H]\bigr)
\quad\text{subject to the displayed kinetic system},
\]
where \(f[H]\) denotes the solution corresponding to \(H\), if such a solution is defined, and \(J\) is intended to measure stability.

This is an incomplete problem specification. The source does not define instability, stabilization, the topology or size of the perturbation, the solution class, the admissible controls, or the time interval. It also does not specify whether a single field must work for every sufficiently small perturbation or whether the field may depend on the perturbation, nor whether feedback is allowed. Its example for \(J\) ends in a truncated electric-field norm expression, so neither the norm nor its temporal evaluation can be fixed. Consequently the displayed minimization is a formal target, not an assertion that a minimizer exists, and no fully quantified stabilization or optimality claim can be extracted without further specifications.
\subsection{Short English statement}
The problem asks whether an external electric field can stabilize an unstable kinetic equilibrium after a small perturbation, and how to choose an optimal such field through PDE-constrained optimization. The supplied equation is retained, but the domains, boundary conditions, control class, meaning of stability, and objective functional require clarification before a precise existence or optimality question can be stated.
\subsection{Sources}
\begin{itemize}
\item Supplied local UnsolvedMath registry record, permanent ID \(30005627\), titled \emph{External-Field Stabilization of Unstable Kinetic Equilibria}; exact source problem code \texttt{catalogue entry}. The supplied record is the local provenance for this draft.
\item Supplied source citation: \emph{Classical and Quantum Mechanical Models of Many-Particle Systems} (2024), Oberwolfach Reports, pp.\ 2179--2246, DOI \texttt{10.4171/owr/2023/38}. Original OWR report: \url{https://doi.org/10.4171/owr/2023/38}. The registry assigns proposed year 2023; the citation's 2024 date and the DOI's 2023 component are preserved as supplied.
\item The extracted original text contains the heading \emph{Two case studies of PDE-constrained optimization within kinetic framework} and the name Qin Li. This is contextual attribution from the excerpt; the registry supplies no proposer attribution.
\item Supplied literature source: Einkemmer--Li--Mouhot--Yue, Vlasov--Poisson control, \url{https://doi.org/10.1016/j.jcp.2025.113904}.
\item Literature assessment dated exactly \texttt{2026-08-22}: the supplied curation labels the problem partially solved and reports explicit external fields suppressing instabilities for known equilibria and perturbations, while describing the general extracted PDE-constrained stabilization problem as unresolved. This is supplied dated curation, not an independently verified literature assessment.
\item The registry labels the statement \texttt{corrected\_verified} and says it was reconstructed and checked against the official Oberwolfach report. These are supplied verification claims; this draft does not independently verify them or claim to have visited either source.
\end{itemize}
\end{document}
